\section{The SCR-ASC System: Mathematical Modelling and Parametrization}

This section develops an identifiable mathematical model for the catalytic $\nox$ reduction dynamics in the SCR-ASC
system. Constraints on sensor placement and sampling rates limit the available measurements to inlet and outlet $\nox$
concentrations and urea dosing sampled at periods large relative to the residence time of the reacting flow. Rather than
modelling the internal reacting-flow dynamics and then discretizing, we model the time evolution of these measurement
signals directly, treating the SCR-ASC chamber as a control volume of known dimensions and accounting for the underlying
reaction kinetics that generate the measurements. Because the discretization follows from the measurement process
itself, the resulting model is natively discrete-time and directly usable for system identification, state estimation,
fault detection, and control design. Assumptions, labeled $A_i$, are introduced and justified at the point of use
throughout the derivation.

\begin{figure}[!ht]
    \centering
    \includegraphics[width=\figWidth]{./figs/2-model/SCR-ASC_ModelReduced.png}
    \caption{SCR-ASC reduced-order model}
    \label{fig::scr_asc_reduced}
\end{figure}

% ===========================================================
\subsection{Reactions and Rate Equations}
%TODO: Add references to the claims
The SCR-ASC system involves urea decomposition, ammonia adsorption/desorption, selective catalytic reduction of $\nox$,
and ammonia oxidation (AMOX) on both SCR and ASC surfaces. The following assumptions reduce the full reaction network to
a tractable form.

\itbf{$A_1$:} All $\nox$ in the exhaust is $\nitrox$, since commercially available $\nox$ sensors cannot differentiate
between $\nitrox$ and $\notwo$ \cite{nova2014urea}.

\itbf{$A_2$:} Three dominant reactions are retained: the standard SCR reaction, ammonia oxidation (AMOX), and ammonia
adsorption/desorption. Fast and slow SCR reactions are omitted, as the exhaust flow rate ensures they do not
significantly affect tailpipe concentrations.

\itbf{$A_3$:} The AMOX reactions on both SCR and ASC surfaces are lumped into a single oxidation reaction with $100\%$
$\nitrogen$ selectivity (valid for $T > 225\,°\text{C}$).

\itbf{$A_4$:} The Eley-Rideal mechanism governs the standard SCR reaction: gaseous $\nox$ reacts with surface-adsorbed
$\ammonia$.

\itbf{$A_5$:} Mass transfer is neglected; the catalyst kinetics are reaction-controlled, as the standard SCR reaction
rate exceeds the exhaust fluid flow rate.

\itbf{$A_6$:} The SCR-ASC chamber is modelled as a plug-flow reactor; reaction rates within each residence time are
governed by inlet gas-phase concentrations.
%===
\begin{enumerate}
    \item Standard SCR reaction:
          \begin{align}
              4 \ammonia^{ads} + 4 \nitrox + \oxygen & \xrightarrow[]{k_{scr}} 4 \nitrogen + 6 \water \label{react::std_scr}
          \end{align}
    \item Ammonia Oxidation:
          \begin{align}
              4 \ammonia^{ads} + 3 \oxygen & \xrightarrow[]{k_{oxi}} 2 \nitrogen + 6 \water \label{react::amox}
          \end{align}
    \item Ammonia Adsorption/Desorption:
          \begin{align}
              \ammonia + \Theta_{free} & \xrightleftharpoons[k_{des}]{k_{ads}} \ammonia^{ads}
              \label{react::ads}
          \end{align}
\end{enumerate}
%===
The surface $\ammonia$ rate follows directly from these three reactions. Adsorption fills sites proportional to
free-site concentration $(\Gamma - \con{\ammonia}^{ads})$, while desorption, SCR, and oxidation each consume adsorbed
$\ammonia$:
\begin{multline}
    \dot{\con{\ammonia}}^{ads} = \underbrace{k_{ads}\,\con{\ammonia}^{in}\,\lr{\Gamma - \con{\ammonia}^{ads}}}_{\text{adsorption}} \\
    - \underbrace{\lr{k_{des} + k_{scr}\,\con{\nox}^{in} + k_{oxi}}\,\con{\ammonia}^{ads}}_{\text{desorption + SCR + oxidation}}
    \label{eqn::nh3_ads_rate}
\end{multline}
where $\Gamma$ and $\con{\ammonia}^{ads}$ are surface concentrations ($mol/m^2$), while the remaining quantities are
volumetric concentrations ($mol/m^3$). The conversion between surface and volumetric domains is given by the
surface-to-volume ratio $k_{s2v} = A_{scr}/V$ ($m^{-1}$). The volumetric $\nox$ consumption rate involves only the SCR
reaction:
\begin{align}
    \dot{\con{\nox}}^{scr} &= -k_{s2v}\,k_{scr}\,\con{\ammonia}^{ads}\,\con{\nox}^{in}
    \label{eqn::nox_rate}
\end{align}
These two rate equations are coupled through $\con{\ammonia}^{ads}$. The surface concentration drives $\nox$
reduction but is itself unobservable from tailpipe measurements. Eliminating this hidden state and expressing the $\nox$
dynamics purely in terms of measurable signals is the central objective of the derivation that follows.
% ===========================================================
% §2.2 — Physical Property Models
% Sources: Thesis/secs/2-data/2-data_units.tex (density)
%          Thesis/secs/2-data/1-res_time.tex (residence time)
%          Thesis/secs/3-governing_eqns/1-reactions.tex (Arrhenius)
%          Thesis/secs/4-NOx_mdl/1-subs/parametrization.tex (Chebyshev, A3–A6)
%          Thesis/secs/3-governing_eqns/5-urea_dosing_proc_dyn.tex (urea dosing)
% ===========================================================
\subsection{Physical Property Models}
The governing equations developed in this section are parametrized using approximate models for the physical properties
of the system. These models bridge the gap between the first-principles rate equations and the identifiable parametric
form presented in subsequent subsections.

% --- §2.2.1 Exhaust Gas Density ---
\subsubsection{Exhaust Gas Density}
The density of the exhaust flow is approximated as the density of air at the exhaust temperature and ambient atmospheric
pressure, using the ideal gas law:
\begin{align}
    \rho &= \frac{PM}{R T} = \frac{\mu}{T}
    \label{eqn::density_model}
\end{align}
where $P = 101.325 \, kPa$ is the ambient pressure, $M = 28.97 \, g/mol$ is the mean molecular weight, and $R = 8.314
\, J/(mol \cdot K)$ is the universal gas constant. At the typical operating temperature ($T_0 \approx 250 \lx{^o}{C}
\approx 523 \, K$), a $\pm 100 \lx{^o}{C}$ variation yields less than $10\%$ change in density, which is negligible
compared to the change in mass flow rate. Similarly, the sensitivity of density to $\nox$ concentration is negligible
because the mole fraction of $\nox$ in exhaust is at the parts-per-million (ppm) level. Thus, density is assumed
constant:
\begin{align}
    \rho &\approx \rho_0 = 1.2 \times 10^{-3} \, g/cm^3 \quad \text{at} \quad 250 \lx{^o}{C}
    \label{eqn::const_density}
\end{align}

% --- §2.2.2 Residence Time ---
\subsubsection{Residence Time Model}
The residence time of the reacting flow through the SCR-ASC chamber is defined as:
\begin{align}
    \tau &= \frac{\rho V_{scr}}{F}
    \label{eqn::res_time_def}
\end{align}
where $V_{scr}$ is the volume of the SCR-ASC chamber and $F$ is the mass flow rate. Using the constant-density
assumption~(\ref{eqn::const_density}):
\begin{align}
    \tau(k) &= \frac{\rho_0 V_{scr}}{F(k)} = \frac{\tau_0}{F(k)}
    \label{eqn::residence_time_mdl}
\end{align}
where $\tau_0 = \rho_0 V_{scr}$ is a constant. The mode of the residence time distribution for the available test-cell
and truck data is approximately $0.07 \, s$, which is much shorter than the sampling intervals ($t_s^{test} = 0.2 \, s$,
$t_s^{truck} = 1 \, s$). Thus, multiple residence times occur within each sampling period ($n = t_s / \tau \gg 1$),
motivating the telescopic summation approach adopted in this work.

% --- §2.2.3 Rate Constant Models (Arrhenius → Chebyshev) ---
\subsubsection{Rate Constant Temperature Models}
The rate constants of the reactions follow the Arrhenius equation:
\begin{align}
    k &= A \exp\left(-\frac{E}{RT}\right)
\end{align}
where $A$ is the pre-exponential factor and $E$ is the activation energy. Using Taylor series expansion about a nominal
temperature $T_0$, the rate constant can be approximated as a polynomial in temperature:
\begin{align}
    k(T) &\approx mT + c     && \text{(linear, valid within } \pm 50 \lx{^o}{C}\text{)}
        \label{eqn::k_linear} \\
    k(T) &\approx qT^2 + mT + c && \text{(quadratic, wider range)}
        \label{eqn::k_quadratic}
\end{align}
For numerical stability, the temperature range $[T_{min}, T_{max}]$ is mapped to $[-1, 1]$ and Chebyshev polynomial
basis functions are used:
\begin{align}
    \phi_1 (k) &= \bm{ \frac{T-T_0}{T_r} & 1}
    \label{eqn::phi_1} \\
    \phi_2 (k) &= \bm{2\lr{\frac{T-T_0}{T_r}}^2-1 & \frac{T-T_0}{T_r} & 1}
    \label{eqn::phi_2}
\end{align}
where $T_0 = (T_{max} + T_{min})/2$ and $T_r = (T_{max} - T_{min})/2$. The first-order basis $\phi_1$ is used for
terms containing individual rate constants, which vary nearly linearly over the operating range. The second-order basis
$\phi_2$ is used for terms containing products of rate constants and $\Gamma$ (e.g., $k_{scr} \cdot k_{ads} \cdot
\Gamma$), where the combined temperature dependence includes an inflection point.

% --- §2.2.4 Urea Dosing Model ---
\subsubsection{Urea Dosing Model}
The primary control input to the system is the injection of urea (commercially known as AdBlue, a $32.5\%$ aqueous urea solution). The aggregate decomposition reaction is:
\begin{align*}
    \urea \, (liquid) &\longrightarrow 2 \ammonia + \cotwo + x \water
\end{align*}
Since the concentration of urea in the aqueous solution is constant, the rate of ammonia production is directly proportional to the urea injection rate. The inlet volumetric concentration of ammonia is given by:
\begin{align}
    \con{\ammonia}^{in}(k) &= \nu_u \frac{u_{inj}(k)}{F(k)}
    \label{eqn::urea_dosing}
\end{align}
where $\nu_u = 2 r_u t_s \rho_0$ is a lumped constant independent of temperature (since urea is preheated before
injection), $u_{inj}$ is the urea injection rate, and $F$ is the exhaust mass flow rate. This reciprocal dependence on
$F$ breaks down at zero flow rate, which is physically consistent.

% ===========================================================
\subsection{$\nox$ Reduction Model: Eliminating the Unobservable $\sigma$ \label{sec::nox_io_model}}

The internal state $\sigma(k)$ is not directly measurable. To obtain an identifiable input-output model, $\sigma$ is
eliminated using the $\nox$ output equation~(\ref{eqn::nox_output}). Define the \textit{$\nox$ reduction per sample}:
\begin{align}
    \eta(k) &= u_1(k{-}1) - x_1(k)
    \label{eqn::eta_def}
\end{align}
where $u_1(k) = \con{\nox}^{in}(k)$ and $x_1(k) = \con{\nox}^{out}(k)$. From~(\ref{eqn::nox_output}):
\begin{align}
    \eta(k{+}1) &= \tau(k)\, k_{s2v}\, k_{scr}(k)\, u_1(k)\, \sigma(k)
\end{align}
Hence the average surface concentration can be expressed as:
\begin{align}
    \sigma(k) &= \frac{\eta(k{+}1)}{\tau(k)\, k_{s2v}\, k_{scr}(k)\, u_1(k)}
    \label{eqn::sigma_elim}
\end{align}

% --- Unsaturated recursion ---
\subsubsection{Unsaturated $\eta$-Dynamics}
Rewrite the unbounded surface dynamics~(\ref{eqn::sigma_dynamics}) compactly as:
\begin{align}
    \sigma(k) &= \sigma(k{-}1)\,\gamma_{proc}(k{-}1) + \Gamma\, t_s\, k_{ads}(k{-}1)\, \con{\ammonia}^{in}(k{-}1)
    \label{eqn::sigma_compact}
\end{align}
where
\begin{align}
    \gamma_{proc}(k) &= 1 - t_s\lr{k_{ads}(k)\,\con{\ammonia}^{in}(k) + k_{scr}(k)\,u_1(k) + k_{od}(k)}
    \label{eqn::gamma_proc}
\end{align}
collects all terms that act on the current surface state. Substituting~(\ref{eqn::sigma_elim}) at times $k$ and
$k{-}1$ into~(\ref{eqn::sigma_compact}) eliminates $\sigma$ entirely:
\begin{align}
    f_\sigma(k) =& \;\eta(k) \lrb{\frac{\tau(k)}{\tau(k{-}1)}}
                            \lrb{\frac{u_1(k)}{u_1(k{-}1)}}
                            \lrb{\frac{k_{scr}(k)}{k_{scr}(k{-}1)}}
                            \gamma_{proc}(k{-}1) \notag \\
        &+ t_s\, k_{s2v}\, \Gamma\, \tau(k)\, u_1(k)\, \con{\ammonia}^{in}(k{-}1)\,
            k_{scr}(k)\, k_{ads}(k{-}1)
    \label{eqn::f_sigma}
\end{align}
This is the unsaturated $\eta$-recursion, expressed entirely in terms of $\eta$ (observable) and the measurable inputs
$u_1$, $u_2$, $F$, and $T$.

% --- Switching in η space ---
\subsubsection{Switched Dynamics in $\eta$ Space}
Since $\sigma \propto \eta$ via~(\ref{eqn::sigma_elim}) with a positive proportionality constant, the physical
constraint $0 \leq \sigma \leq \Gamma$ from~(\ref{eqn::g_sat}) maps directly to bounds on $\eta$:
\begin{align}
    0 \;\leq\; \eta(k{+}1) \;\leq\;
        \underbrace{\tau(k)\, k_{s2v}\, k_{scr}(k)\, u_1(k)\, \Gamma}_{\eta_{\sat}(k)}
    \label{eqn::eta_bounds}
\end{align}
where $\eta_{\sat}(k)$ is the $\nox$ reduction when the catalyst is fully saturated ($\sigma = \Gamma$). The switched
dynamics are:
\begin{align}
    \eta(k{+}1) = \begin{cases}
        f_\sigma(k)       & \text{if } 0 \leq f_\sigma(k) \leq \eta_{\sat}(k) \\[4pt]
        \eta_{\sat}(k)    & \text{if } f_\sigma(k) > \eta_{\sat}(k) \\[4pt]
        0                 & \text{if } f_\sigma(k) < 0
    \end{cases}
    \label{eqn::eta_switched}
\end{align}
The switching is entirely in terms of $\eta$ and the measurable inputs ($u_1$, $F$, $T$), with no dependence on the
unobservable $\sigma$.

% ===========================================================
\subsection{Parametrization and Switched Model \label{sec::parametrization}}

The switched $\eta$-dynamics~(\ref{eqn::eta_switched}) are exact but contain individual rate constants and ratio
prefactors. This section parametrizes all three cases ($f_\sigma$, $\eta_{\sat}$, and $0$) to obtain an identifiable
model that is linear in parameters.

% --- Chebyshev basis ---
\subsubsection{Temperature Model for Rate Constants}

\itbf{$A_8$:} The rate constants vary polynomially with temperature over the operating range. This is justified by a
Taylor expansion of the Arrhenius form $k = A\exp(-E/RT)$, which yields a polynomial in $T$ for small perturbations
about a reference temperature.

For numerical stability, the operating range $[T_{min}, T_{max}]$ is mapped to $[-1,1]$ using Chebyshev polynomials.
Let $T_0 = (T_{max}+T_{min})/2$ and $T_r = (T_{max}-T_{min})/2$:
\begin{align}
    \phi_1 (k) &= \bm{ \frac{T-T_0}{T_r} & 1}
    \label{eqn::phi_1} \\
    \phi_2 (k) &= \bm{2\lr{\frac{T-T_0}{T_r}}^2-1 & \frac{T-T_0}{T_r} & 1}
    \label{eqn::phi_2}
\end{align}
First-order $\phi_1$ is used for individual rate constants (nearly linear in $T$). Second-order $\phi_2$ is used for
products of rate constants and $\Gamma$, which exhibit an inflection point.

% --- Parametrize f_sigma ---
\subsubsection{Unsaturated Dynamics: Flow-Scaled Parametric Form}

\itbf{$A_9$:} Temperature changes slowly relative to the sampling period, i.e., $T(k) \approx T(k{-}1)$. This implies
$k_{scr}(k)/k_{scr}(k{-}1) \approx 1$ and allows the product $k_{scr}(k)\, k_{ads}(k{-}1)$ to be treated as a single
combined rate constant $k_{scr/ads}(k{-}1)$.

\itbf{$A_{10}$:} Density changes with temperature are negligible, so $\tau(k) = \tau_0/F(k)$ and
$\tau(k)/\tau(k{-}1) = F(k{-}1)/F(k)$.

Applying $A_9$--$A_{10}$ to the unsaturated recursion~(\ref{eqn::f_sigma}), the ratio prefactor simplifies to
$u_1(k)\, F(k{-}1) / [u_1(k{-}1)\, F(k)]$. To absorb this, define the \textit{flow-scaled $\nox$ reduction}:
\begin{align}
    \bar \eta_F (k) = \frac{F(k{-}1)}{u_1(k{-}1)} \eta(k)
    \label{eqn::eta_bar_def}
\end{align}
Multiplying~(\ref{eqn::f_sigma}) by $F(k)/u_1(k)$ converts $\eta \to \bar\eta_F$ on both sides (the ratio prefactor
becomes unity).

\itbf{$A_{11}$:} The inlet $\ammonia$ concentration follows the urea dosing
model~(\ref{eqn::urea_dosing}): $\con{\ammonia}^{in}(k) = \nu_u\, u_2(k)/F(k)$.

\itbf{$A_{12}$:} The total surface site concentration $\Gamma$ is constant over the operating range (it changes only
with aging).

Expanding $\gamma_{proc}$~(\ref{eqn::gamma_proc}), substituting $A_{11}$--$A_{12}$, and replacing each rate constant
with its Chebyshev basis ($k_{ads} = \phi_1^T\theta_{ads}$, $k_{od} = \phi_1^T\theta_{od}$,
$k_{scr} = \phi_1^T\theta_{scr}$, $k_{scr/ads}\,\Gamma = \phi_2^T\theta_{scr/ads}$):
\begin{align}
    \bar \eta_F (k{+}1) = \bar \eta_F (k) &+ u_{2F}(k{-}1)\,\phi_2(k{-}1)\, \theta_\Gamma
        \notag\\
    &- u_{2F}(k{-}1)\, \bar \eta_F (k)\, \phi_1(k{-}1)\, \theta_{\eta_{ads}}
        \notag\\
    &- \bar \eta_F (k)\, \phi_1(k{-}1)\, \theta_{\eta_{od}}
        \notag\\
    &- u_1(k{-}1)\, \bar \eta_F (k)\, \phi_1(k{-}1)\, \theta_{\eta_{scr}}
    \label{eqn::eta_bar_dynamics}
\end{align}
where $u_{2F}(k) = u_2(k)/F(k)$, and the four parameter groups are:
\begin{inparaenum}[(1)]
    \item $\theta_\Gamma = t_s\, k_{s2v}\, \nu_u\, \Gamma\, \tau_0\, \theta_{scr/ads}$:
        adsorption onto free sites;
    \item $\theta_{\eta_{ads}} = t_s\, \nu_u\, \theta_{ads}$:
        site-blocking by adsorbed $\ammonia$;
    \item $\theta_{\eta_{od}} = t_s\, \theta_{od}$:
        oxidation and desorption loss; and
    \item $\theta_{\eta_{scr}} = t_s\, \theta_{scr}$:
        SCR consumption of adsorbed $\ammonia$.
\end{inparaenum}

% --- Identifiable compact form ---
The model can be written compactly as:
\begin{align}
    \bar \eta_F (k{+}1) &= \bar \eta_F (k) + \phi_{\nox}^T(k)\, \theta_{\nox}
    \label{eqn::eta_compact}
\end{align}
where $\phi_{\nox}(k)$ is the regressor vector (depending only on measured quantities) and
$\theta_{\nox} = \lrb{\theta_{\eta_{ads}}^T \; \theta_{\eta_{od}}^T \; \theta_{\eta_{scr}}^T \;
\theta_\Gamma^T}^T$.

% --- Parametrize eta_sat ---
\subsubsection{Saturated Ceiling}

When $\sigma = \Gamma$, the $\nox$ reduction is $\eta_{\sat}(k) = \tau(k)\, k_{s2v}\, k_{scr}(k)\, u_1(k)\, \Gamma$.
Converting to $\bar\eta_F$ space (multiply by $F(k)/u_1(k)$) and parametrizing $k_{scr}\,\Gamma$ with $\phi_2$
(quadratic, because $k_{scr} \times \Gamma$ has a non-monotone temperature dependence):
\begin{align}
    \bar\eta_{F,\sat}(k) &= \phi_2^T(k)\, \theta_{\sat}
    \label{eqn::eta_sat}
\end{align}
where $\theta_{\sat} = \Gamma\, \tau_0\, k_{s2v}\, \theta_{scr}$ is a separate parameter vector (not shared with the
unsaturated dynamics).

% --- Guard condition ---
\subsubsection{Guard Condition}

The catalyst remains unsaturated when the net adsorption onto free sites is non-negative.
In~(\ref{eqn::eta_bar_dynamics}), the adsorption-related terms are
$u_{2F}\,\phi_2\,\theta_\Gamma - u_{2F}\,\bar\eta_F\,\phi_1\,\theta_{\eta_{ads}}$. Setting this $\geq 0$ and
rearranging into a single inner product:
\begin{align}
    \underbrace{\bm{ \bar \eta_F(k)\, \phi_1^T(k{-}1)
                    & \pmb 0 & \pmb 0 & - \phi_2^T(k{-}1) }}_{\phi_g(k)^T}
                    \theta_{\nox} \leq 0
    \label{eqn::guard_condition}
\end{align}
This guard condition is linear in the model parameters and can be evaluated from the current state and inputs without
knowledge of $\sigma$. The threshold depends on the state $\bar\eta_F(k)$ itself, making the model a
\textit{self-excited threshold nonlinear autoregressive with exogenous input} (SETNARX) structure.

% --- Switched model ---
\subsubsection{Switched Model}

Using the guard condition~(\ref{eqn::guard_condition}) and the saturation ceiling~(\ref{eqn::eta_sat}), let
$G(k) = \phi_g^T(k)\,\theta_{\nox}$. The switched dynamics are:
\begin{align}
    \bar\eta_F(k{+}1) = \begin{cases}
        \bar\eta_F(k) + \phi_{\nox}^T\theta_{\nox}
            & G(k) \leq 0,\; \bar\eta_F > 0 \\[4pt]
        \phi_2^T(k)\, \theta_{\sat}
            & G(k) > 0 \\[4pt]
        0
            & \bar\eta_F = 0,\; G(k) \leq 0
    \end{cases}
    \label{eqn::switched_model}
\end{align}

\begin{figure}[!ht]
    \centering
    \includegraphics[width=\figWidth]{./figs/2-model/hybrid_model.png}
    \caption{Switched $\nox$ model state diagram}
    \label{fig::hybrid_model}
\end{figure}


% ===========================================================
